\documentclass[11pt, a4paper]{article}

\input{bfw-style.tex}

\title{\vspace*{-1.5cm}\Huge\bfseries\color{bfwPrimaryDark}Aufgabenheft Session~5\\[0.4em]
  \large\color{bfwInkSoft}\textnormal{Geschäftsprozesse, Produktionsfaktoren \& Güterarten --- Übungen im IHK-Klausurformat}}
\author{\textsf{IT Lernfeld 1 --- Das Unternehmen und die eigene Rolle im Betrieb beschreiben}}
\date{Selbstlernmaterial}

\begin{document}
\maketitle
\thispagestyle{empty}

\begin{merksatz}[title={\bfseries So nutzen Sie dieses Aufgabenheft}]
Bearbeiten Sie alle Aufgaben zunächst \emph{ohne} in die Lösungen zu schauen. Schreiben
Sie Ihre Antworten in vollständigen Sätzen --- die IHK-Prüfung verlangt formulierte
Begründungen, nicht bloße Stichworte. Beachten Sie die \emph{Operatoren} (nennen,
erläutern, beschreiben, begründen, beurteilen) --- sie geben den geforderten
Antwort-Tiefegrad vor. Erst danach öffnen Sie den Lösungsteil ab Seite~\pageref{loesungen}
und vergleichen.
\end{merksatz}

\bigskip

\textbf{Kontext:} Sie sind Auszubildende bei der \textbf{JIKU IT-Solutions GmbH} in Hamburg.
Ihre Ausbilderin Frau Meyerhoff hat Sie gebeten, die wirtschaftlichen Grundlagen des
Unternehmens zu analysieren --- von der Wertschöpfung über die Geschäftsprozesse bis hin
zu den eingesetzten Gütern.

\bigskip

\noindent
\begin{tabular}{|l|l|l|l|l|l|l|l|}
\hline
\textbf{Aufgabe} & 1 & 2 & 3 & 4 & 5 & 6 & \textbf{Gesamt}\\
\hline
\textbf{Punkte} & /10 & /8 & /10 & /10 & /8 & /10 & /56\\
\hline
\end{tabular}

\tableofcontents
\vspace{1em}
\hrule

\section{Aufgaben}

\subsection*{Aufgabe~1 --- Wertschöpfungsberechnung (10 Punkte)}

JIKU IT-Solutions hat im vergangenen Geschäftsjahr folgende Zahlen erzielt:
Umsatzerlöse (netto) 6,5~Mio.\,€, Vorleistungen (eingekaufte IT-Komponenten und
Fremddienstleistungen) 2,8~Mio.\,€, Abschreibungen auf Sachanlagen 0,4~Mio.\,€.
Die Wertschöpfung wurde zu 38\,\% an Mitarbeiter, zu 22\,\% an Fremdkapitalgeber,
zu 18\,\% an den Staat und zum Rest an Eigenkapitalgeber ausgeschüttet.
JIKU beschäftigt 40 Mitarbeiter.

\begin{enumerate}[label=(\alph*)]
  \item \textbf{Berechnen} Sie die Wertschöpfung von JIKU IT-Solutions. \hfill (3~P.)

  \vspace{5cm}

  \item \textbf{Berechnen} Sie die Wertschöpfung pro Mitarbeiter (W/M). \hfill (2~P.)

  \vspace{4cm}

  \item \textbf{Berechnen} Sie die Wertschöpfungstiefe. \hfill (2~P.)

  \vspace{4cm}

  \item \textbf{Ermitteln} Sie den absoluten Betrag, der an die Eigenkapitalgeber
  ausgeschüttet wird. \hfill (3~P.)

  \vspace{4.5cm}
\end{enumerate}

\subsection*{Aufgabe~2 --- Wertschöpfungskette (8 Punkte)}

\begin{enumerate}[label=(\alph*)]
  \item \textbf{Beschreiben} Sie, was unter der Wertschöpfungskette eines Unternehmens
  zu verstehen ist, und \textbf{erläutern} Sie, welche Tätigkeiten ausdrücklich nicht
  dazuzählen. \hfill (3~P.)

  \vspace{5cm}

  \item \textbf{Ordnen} Sie die folgenden Tätigkeiten bei JIKU der korrekten
  Wertschöpfungskettenposition zu (Beschaffung, Marketing, Verkauf, Dienstleistung/Service)
  oder begründen Sie, warum sie nicht zur Wertschöpfungskette gehören. \hfill (5~P.)

  \bigskip

  \noindent
  \begin{tabular}{|p{4.5cm}|p{3.5cm}|p{5cm}|}
  \hline
  \textbf{Tätigkeit} & \textbf{Position} & \textbf{Begründung}\\
  \hline
  Netzwerkinstallation beim Kunden & & \\[0.8cm]
  \hline
  Lohnabrechnung der Mitarbeiter & & \\[0.8cm]
  \hline
  Anfrage bei Lieferanten & & \\[0.8cm]
  \hline
  Angebotserstellung für Neukunden & & \\[0.8cm]
  \hline
  After-Sales-Support per Telefon & & \\[0.8cm]
  \hline
  \end{tabular}
\end{enumerate}

\subsection*{Aufgabe~3 --- Geschäftsprozesse (10 Punkte)}

Frau Meyerhoff erklärt Ihnen, dass bei JIKU verschiedene Prozesstypen unterschieden werden.

\begin{enumerate}[label=(\alph*)]
  \item \textbf{Unterscheiden} Sie Kernprozesse und Supportprozesse hinsichtlich
  Wertschöpfung, Kundenorientierung und Beispielen. \hfill (4~P.)

  \vspace{6cm}

  \item \textbf{Ordnen} Sie folgende JIKU-Aktivitäten dem richtigen Prozesstyp zu
  (Kernprozess, Supportprozess oder Führungsprozess) und \textbf{begründen} Sie kurz. \hfill (6~P.)

  \bigskip

  \noindent
  \begin{tabular}{|p{6.5cm}|p{2.5cm}|p{4.5cm}|}
  \hline
  \textbf{Aktivität} & \textbf{Prozesstyp} & \textbf{Begründung}\\
  \hline
  Netzwerkeinrichtung beim Kunden & & \\[0.8cm]
  \hline
  Jahresplanung durch Geschäftsleitung & & \\[0.8cm]
  \hline
  Pflege der internen JIKU-Server & & \\[0.8cm]
  \hline
  Cloud-Migration für Mittelstandskunden & & \\[0.8cm]
  \hline
  Monatliche Gehaltsabrechnungen & & \\[0.8cm]
  \hline
  Ressourcenplanung für nächstes Quartal & & \\[0.8cm]
  \hline
  \end{tabular}
\end{enumerate}

\subsection*{Aufgabe~4 --- Betriebliche Produktionsfaktoren nach Gutenberg (10 Punkte)}

\begin{enumerate}[label=(\alph*)]
  \item \textbf{Nennen} Sie die drei Elementarfaktoren nach Gutenberg und
  \textbf{erläutern} Sie je einen kurz. \hfill (6~P.)

  \vspace{7cm}

  \item \textbf{Ordnen} Sie folgende Ressourcen bei JIKU dem richtigen Produktionsfaktor
  zu und \textbf{begründen} Sie kurz. \hfill (4~P.)

  \bigskip

  \noindent
  \begin{tabular}{|p{5.5cm}|p{3cm}|p{4.5cm}|}
  \hline
  \textbf{Ressource} & \textbf{Faktor} & \textbf{Begründung}\\
  \hline
  Netzwerkkabel und Festplatten & & \\[0.8cm]
  \hline
  Serverracks im Rechenzentrum & & \\[0.8cm]
  \hline
  IT-Techniker bei Kundeninstallation & & \\[0.8cm]
  \hline
  Standortleiterin, koordiniert Personal & & \\[0.8cm]
  \hline
  \end{tabular}
\end{enumerate}

\subsection*{Aufgabe~5 --- Güterarten (8 Punkte)}

\begin{enumerate}[label=(\alph*)]
  \item \textbf{Erläutern} Sie den Unterschied zwischen freien und wirtschaftlichen Gütern
  sowie zwischen Verbrauchs- und Gebrauchsgütern. Geben Sie je ein Beispiel aus dem
  JIKU-Umfeld. \hfill (4~P.)

  \vspace{6cm}

  \item \textbf{Klassifizieren} Sie eine Softwarelizenz (z.\,B.\ Betriebssystem) nach
  \emph{allen} relevanten Gütermerkmalen (Beschaffenheit, Knappheit, Ort/Zweck,
  Ausschließbarkeit) und \textbf{begründen} Sie jeden Aspekt. \hfill (4~P.)

  \vspace{6cm}
\end{enumerate}

\subsection*{Aufgabe~6 --- Fallaufgabe: Digitale Transformation bei JIKU (10 Punkte)}

JIKU IT-Solutions überlegt, seinen Kunden eine eigene digitale Plattform für IT-Support
und Remote-Monitoring anzubieten, die rund um die Uhr verfügbar ist. Gleichzeitig möchte
JIKU seinen bisherigen papierbasierten Bestellprozess für IT-Komponenten durch ein
E-Procurement-System ersetzen.

\begin{enumerate}[label=(\alph*)]
  \item \textbf{Unterscheiden} Sie, welche der beiden Maßnahmen als „digitale
  Transformation" und welche --- falls überhaupt --- als „digitale Disruption" eingestuft
  werden kann. \textbf{Begründen} Sie Ihre Einschätzung. \hfill (4~P.)

  \vspace{6cm}

  \item \textbf{Benennen} Sie für den Bestellprozess das geeignete integrierte
  Softwaresystem (Kürzel genügt) und \textbf{erläutern} Sie kurz, was dieses System
  leistet. \hfill (2~P.)

  \vspace{4cm}

  \item \textbf{Erklären} Sie, warum die eigene Support-Plattform als Kernprozess
  einzustufen ist, und \textbf{nennen} Sie zwei Güterkategorien, die JIKU über die
  Plattform anbietet. \hfill (4~P.)

  \vspace{5cm}
\end{enumerate}

% =====================================================================
\newpage
\section*{Lösungshinweise für Lehrende}
\label{loesungen}
\addcontentsline{toc}{section}{Lösungshinweise für Lehrende}
% =====================================================================

\begin{merksatz}[title={\bfseries Hinweis zur Nutzung}]
Dieser Abschnitt ist ausschließlich für Lehrende bestimmt. Die Lösungen geben
Musterlösungen und typische Fehlerquellen an. Teilnehmende erhalten diesen Teil
erst nach der eigenständigen Bearbeitung.
\end{merksatz}

\subsection*{Lösung Aufgabe~1}

\begin{enumerate}[label=(\alph*)]
  \item Wertschöpfung = Umsatzerlöse minus Vorleistungen minus Abschreibungen\\
  = 6{,}5~Mio.\,€ $-$ 2{,}8~Mio.\,€ $-$ 0{,}4~Mio.\,€ = \textbf{3{,}3~Mio.\,€}

  \textit{Häufiger Fehler: Personalkosten als Abzugsposten ansetzen. Löhne und Gehälter sind
  Teil der Wertschöpfungsverteilung, nicht der Berechnung.}

  \item W/M = 3{,}3~Mio.\,€ $\div$ 40 Mitarbeiter = \textbf{82.500~€ pro Mitarbeiter}

  \item Wertschöpfungstiefe = 3{,}3~Mio.\,€ $\div$ 6{,}5~Mio.\,€ = \textbf{ca. 0{,}508} (rund 50{,}8\,\%)

  \item Anteil Eigenkapitalgeber = 100\,\% $-$ 38\,\% $-$ 22\,\% $-$ 18\,\% = 22\,\%\\
  Absolutbetrag = 3{,}3~Mio.\,€ × 0{,}22 = \textbf{726.000~€}
\end{enumerate}

\subsection*{Lösung Aufgabe~2}

\begin{enumerate}[label=(\alph*)]
  \item Die Wertschöpfungskette zeigt die kundennah und wertschöpfungsintensiven
  Tätigkeitsbereiche eines Unternehmens. Jedes Unternehmen versucht, seine
  Wertschöpfungskette anzupassen, um eine möglichst hohe Wertschöpfung zu erzielen.
  Verwaltende Tätigkeiten außerhalb der Kette (z.\,B.\ Personalwesen, Rechnungswesen,
  interne IT) zählen ausdrücklich \emph{nicht} zur Wertschöpfungskette.

  \item
  \begin{itemize}
    \item Netzwerkinstallation beim Kunden: \textbf{Dienstleistung} --- kundennah und
    wertschöpfungsintensiv.
    \item Lohnabrechnung der Mitarbeiter: \textbf{nicht in der Wertschöpfungskette} ---
    interne Verwaltungsaufgabe (Supportprozess).
    \item Anfrage bei Lieferanten: \textbf{Beschaffung} --- Einstieg in die Wertschöpfungskette.
    \item Angebotserstellung für Neukunden: \textbf{Verkauf} --- direkte Kundeninteraktion
    zur Leistungsanbahnung.
    \item After-Sales-Support: \textbf{Service} --- letztes Glied der Wertschöpfungskette.
  \end{itemize}
\end{enumerate}

\subsection*{Lösung Aufgabe~3}

\begin{enumerate}[label=(\alph*)]
  \item \textbf{Kernprozesse:} kundennah und wertschöpfungsintensiv; das Unternehmen weist
  Kernkompetenzen nach (z.\,B.\ IT-Systembau, Serviceleistungen). Ohne Kernprozesse keine
  Erlöse. \textbf{Supportprozesse:} nicht direkt wertschöpfend, aber notwendig für die
  Kernprozesse; erzeugen keinen direkten Kundennutzen (z.\,B.\ Buchhaltung, Gehaltsabrechnung).

  \item
  \begin{itemize}
    \item Netzwerkeinrichtung beim Kunden: \textbf{Kernprozess} --- direkter Kundennutzen, wertschöpfungsintensiv.
    \item Jahresplanung durch Geschäftsleitung: \textbf{Führungsprozess} --- koordiniert und steuert das Unternehmen auf Zielebene.
    \item Pflege der internen JIKU-Server: \textbf{Supportprozess} --- interne Infrastruktur ohne direkten Kundennutzen.
    \item Cloud-Migration für Mittelstandskunden: \textbf{Kernprozess} --- kundennah, wertschöpfungsintensiv, Kernkompetenz.
    \item Monatliche Gehaltsabrechnungen: \textbf{Supportprozess} --- interne Verwaltungsaufgabe.
    \item Ressourcenplanung für nächstes Quartal: \textbf{Führungsprozess} --- Planung und Organisation durch Leitung.
  \end{itemize}
  \textit{Hinweis: Gutenberg zeigt die Führungsaufgabe als dispositiven Faktor; auf Prozessebene
  spricht man von Führungsprozessen (nicht im Lehrbuch-Kapitel explizit, jedoch im Quiz behandelt).}
\end{enumerate}

\subsection*{Lösung Aufgabe~4}

\begin{enumerate}[label=(\alph*)]
  \item Drei Elementarfaktoren nach Gutenberg:
  \begin{itemize}
    \item \textbf{Werkstoffe:} Alle Güter, die zur Erstellung eines Produktes notwendig sind:
    Rohstoffe (Hauptbestandteil), Hilfsstoffe (Nebenbestandteil), Betriebsstoffe (nicht
    direkt zurechenbar), Fertigteile.
    \item \textbf{Betriebsmittel:} Technische Ausstattungen, die im Produktionsprozess
    verwendet werden: Betriebsgebäude, Maschinen, Systeme, Fahrzeuge.
    \item \textbf{Ausführende Arbeit:} Menschliche Arbeitsleistung \emph{ohne} Leitungs-,
    Planungs-, Organisations- oder Kontrollfunktion.
  \end{itemize}
  \textit{Ergänzend: Dispositiver Faktor = Geschäftsleitung mit Leitung, Planung, Organisation,
  Kontrolle; Information als moderner Ergänzungsfaktor.}

  \item
  \begin{itemize}
    \item Netzwerkkabel und Festplatten: \textbf{Werkstoffe} (Hilfs- bzw. Einbauteile, zur Leistungserstellung notwendig).
    \item Serverracks im Rechenzentrum: \textbf{Betriebsmittel} (technische Ausstattung).
    \item IT-Techniker bei Kundeninstallation: \textbf{Ausführende Arbeit} (operativ, ohne Leitungsfunktion).
    \item Standortleiterin, koordiniert Personal: \textbf{Dispositiver Faktor} (Leitung und Organisation).
  \end{itemize}
\end{enumerate}

\subsection*{Lösung Aufgabe~5}

\begin{enumerate}[label=(\alph*)]
  \item \textbf{Freie Güter} sind in der Natur unbegrenzt vorhanden und haben keinen Preis
  (z.\,B.\ Atemluft). \textbf{Wirtschaftliche Güter} sind knapp und haben einen Preis
  (z.\,B.\ Netzwerkkabel, Strom für das Rechenzentrum). \textbf{Verbrauchsgüter} werden
  bei der Nutzung einmalig verbraucht (z.\,B.\ Druckerpapier, Stromverbrauch der Server).
  \textbf{Gebrauchsgüter} können wiederholt genutzt werden (z.\,B.\ Serverrack, Arbeitsplatz\-rechner).

  \item Softwarelizenz:
  \begin{itemize}
    \item \textbf{Beschaffenheit:} immateriell (kein physisches Produkt; es ist ein Nutzungsrecht).
    \item \textbf{Knappheit:} wirtschaftliches Gut (kostenpflichtig, begrenzte Lizenzen verfügbar).
    \item \textbf{Ort/Zweck:} Investitionsgut (zur betrieblichen Nutzung) und Gebrauchsgut
    (wiederholt über die Lizenzlaufzeit nutzbar).
    \item \textbf{Ausschließbarkeit:} privates Gut (ohne Lizenzschlüssel kein Zugang;
    Ausschlussmöglichkeit durch DRM und Lizenzverwaltung).
  \end{itemize}
\end{enumerate}

\subsection*{Lösung Aufgabe~6}

\begin{enumerate}[label=(\alph*)]
  \item \textbf{E-Procurement-Einführung:} digitale Transformation --- ein bestehender
  Prozess (Bestellung) wird von analog auf digital umgestellt; die grundsätzliche Art der
  Beschaffung bleibt gleich. \textbf{Support-Plattform:} ebenfalls digitale Transformation
  (bestehende Serviceleistung über neuen Kanal 24/7 verfügbar machen). Echte Disruption
  würde bedeuten, das Kundenproblem auf eine komplett andere, neuartige Weise zu lösen ---
  z.\,B.\ eine KI-gestützte Plattform, die selbst IT-Fehler behebt, ohne dass JIKU manuell
  eingreift (ganz neues Geschäftsmodell).

  \textit{Hinweis: Gut begründete Einschätzungen, die die Plattform als disruptions\-nahe einstufen,
  sind vertretbar, wenn der Wechsel vom reaktiven Service zum proaktiven KI-System beschrieben wird.}

  \item \textbf{E-P} (E-Procurement): elektronische Beschaffung von Gütern über das Internet
  und spezielle Beschaffungsplattformen; ersetzt papierbasierte Bestellprozesse und verbessert
  Transparenz und Geschwindigkeit.

  \item Die Support-Plattform ist ein \textbf{Kernprozess}, weil sie unmittelbaren Kundennutzen
  stiftet (direkter Kundenkontakt, wertschöpfungsintensiv) und zu den Kernkompetenzen von JIKU
  als IT-Systemhaus gehört. Güterkategorien (Auswahl): \textbf{Dienstleistungen} (Remote-Support,
  Monitoring) und \textbf{immateriell-wirtschaftliche Güter / Rechte} (Zugang zur Plattform als
  kostenpflichtiger Dienst).
\end{enumerate}

\end{document}
